% Transcription littérale des photographies de photos/02.
% Les abréviations et l'ordre des raisonnements de la copie sont conservés.

\providecommand{\Sp}{\operatorname{Sp}}

\exercise{1}{}

\begin{statementbox}
Soit $V$ un $\C$ espace vectoriel de dimension $d \in \N^*$.
\begin{enumerate}[label=Q\arabic*)]
  \item Soit $f \in \Lin(V)$, mq $(\rg(f^n))_{n \in \N}$ est convergente. On note $r(f)$ sa limite.
  \item Soient $(f,g) \in \Lin(V)^2$ qui commutent, mq $r(f+g) \leq r(f) + r(g)$. Est-ce toujours le cas si $f$ et $g$ ne commutent pas ?
  \item Exprimer $r(f)$ en fonction de la multiplicité $m(0)$ de $0$ en tant que racine de $\chi_f$.
\end{enumerate}
\end{statementbox}

\begin{solution}
\partanswer{Q1)}
On a $\forall n \in \N$, $\Ima(f^{n+1}) \subset \Ima(f^n)$
donc $\rg(f^{n+1}) \leq \rg(f^n)$.

La suite des $(\rg(f^n))_{n \in \N}$ est une suite d'entiers naturels décroissante donc stationnaire, donc convergente.

Rem : Si $f$ est bijective, $\forall n \in \N$, $\rg(f^n) = d$ et $r(f) = d$.

\partanswer{Q2)}
Soient $(n,m) \in \N^2$, $\forall p \geq m$, $\Ima(f^p) = \Ima(f^n)$ et si $p \geq m$ $\Ima(g^p) = \Ima(g^m)$.

Soit $k \geq n+m$,
\[
(f+g)^k = \sum_{l=0}^k \binom{k}{l} f^l g^{k-l} \quad \text{car } f \text{ et } g \text{ commutent.}
\]
Soit $x \in V$, si $l \geq n$, $f^l \circ g^{k-l}(x) \in \Ima(f^n)$.
Si $l \leq n$, $k-l \geq n+m-l \geq m$ et
\[
f^l \circ g^{k-l}(x) = g^{k-l} \circ f^l(x) \in \Ima(g^m).
\]
Donc, $\Ima((f+g)^k) \subset \Ima(f^n) + \Ima(g^m)$
soit
\begin{align*}
\rg((f+g)^k) &\leq \dim(\Ima(f^n) + \Ima(g^m)) \\
&\leq \underbrace{\dim \Ima(f^n)}_{r(f)} + \underbrace{\dim \Ima(g^m)}_{r(g)}
\end{align*}
$\forall k \geq n+m$, $\rg((f+g)^k) \leq \rg(f) + \rg(g)$.

Si $f$ et $g$ ne commutent pas,
$f \in \Lin(\C^2)$ c.a. à $\begin{pmatrix} 0 & 1 \\ 0 & 0 \end{pmatrix}$ et $g$ c.a. à $\begin{pmatrix} 0 & 0 \\ 1 & 0 \end{pmatrix}$.
On a $(f+g)$ c.a. à $\begin{pmatrix} 0 & 1 \\ 1 & 0 \end{pmatrix} \in \GL_2(\C)$, donc $r(f+g) = 2$.
Or, $r(f) = r(g) = 0$ car $f$ et $g$ nilpotentes.

\partanswer{c)}
Si $m(0) = 0$, $0$ n'est pas racine de $\chi_f$ donc $0 \notin \Sp(f)$.
Donc, $f$ est bijective et $r(f) = d$.

Si $m(0) = d$, d'après le T.C.H, on a $\chi_f = X^d$ et $f^d = 0$ donc $f$ est nilpotente, soit $r(f) = 0$.

Sinon, on peut écrire $\chi_f = X^{m(0)} Q$ où $Q(0) \neq 0$.
Aussi, $X^{m(0)} \land Q = 1$. D'après le théorème de C-H, $\chi_f$ est un polynôme annulateur de $f$, et d'après le lemme des noyaux, on peut écrire
\[
V = \underbrace{\Ker(f^{m(0)})}_{V_1} \oplus \underbrace{\Ker(Q(f))}_{V_2}
\]
$V_1$ et $V_2$ sont stables par $f$.
Soit $B_1$ une base de $V_1$, $B_2$ une base de $V_2$, $B = B_1 \cup B_2$.
On a $\Mat_B(f) = \begin{pmatrix} A_1 & 0 \\ 0 & A_2 \end{pmatrix}$ avec $A_1 = \Mat_{B_1}(f_{|V_1})$, $A_2 = \Mat_{B_2}(f_{|V_2})$.

Si $f_{|V_2}$ n'était pas injective, $\exists x \in V_2 \setminus \{0\} / f(x) = 0$
et $x \in V_2 \cap \Ker(f) \subset \{0\}$ ce qui ne se peut.
($\Ker(f) \subset \Ker(f^{m(0)})$)

Donc, $f_{|V_2}$ est bijective.
Par définition de $V_1$, $f_{|V_1}^{m(0)} = 0$ donc $A_1^{m(0)} = 0$.
et $\forall k > m(0)$, $\Mat_B(f)^k = \begin{pmatrix} 0 & 0 \\ 0 & A_2^k \end{pmatrix}$ où $A_2$ inversible.
Ainsi, $\rg(f^k) = \dim V_2 = d - \dim V_1$.

\begin{lemma*}
$\dim V_1 = m(0)$.
\end{lemma*}
Dém : Soit $\lambda \in \Sp(f_{|V_1})$. Comme $f_{|V_1}^{m(0)} = 0$, on a $\lambda = 0$ donc $\Sp(f_{|V_1}) = \{0\}$ donc $\chi_{f_{|V_1}} = X^{\dim V_1}$.
Or, $\chi_f = \chi_{f_{|V_1}} \times \chi_{f_{|V_2}}$ et $0$ n'est pas racine de $\chi_{f_{|V_2}}$.
Donc, $\dim V_1 = m(0)$.

Finalement, $r(f) = d - m(0)$.
\end{solution}

\exercise{8}{}

\begin{statementbox}
Soit $V$ un $\C$-ev de dimension finie, $u \in \Lin(V)$ est cyclique ssi $\exists x, (x, \dots, u^{n-1}(x))$ est une base de $V$.
Mq $u$ est cyclique ssi les sous-espaces propres sont de dim $1$.
\end{statementbox}

\begin{solution}
Supposons $u$ cyclique : Soit $x \in V$ tel que $\mathcal B = (x, u(x), \dots, u^{n-1}(x))$ soit une base de $V$.
\[
\Mat_{\mathcal B}(u) =
\begin{pmatrix}
0 & & & \alpha_0 \\
1 & \ddots & & \vdots \\
& \ddots & 0 & \vdots \\
0 & & 1 & \alpha_{n-1}
\end{pmatrix} = A
\]
Soit $\lambda \in \C$ :
\[
A - \lambda I_n =
\begin{pmatrix}
-\lambda & & & \alpha_0 \\
1 & \ddots & & \vdots \\
& \ddots & -\lambda & \vdots \\
0 & & 1 & \alpha_{n-1} - \lambda
\end{pmatrix}
\]
Alors les $n-1$ premières colonnes sont libres car échelonnées.
Donc $\dim \Ker(u - \lambda \Id) \leq 1$. Si $\lambda$ est valeur propre, nécessairement, $\dim \Ker(u - \lambda \Id) = 1$.

Réciproquement : s'il existe $\lambda \in \C$ et $r \in \N / \pi_u = (X-\lambda)^r$.
Soit $\alpha \in \Sp(u)$, nécessairement $\alpha = \lambda$ [$\pi_u$ annule $\alpha$].
D'autre part, $u$ possède nécessairement une valeur propre car $u \in \Lin(V)$ et $V$ est un $\C$-ev de dimension finie non nulle. Donc $\Sp(u) = \{\lambda\}$.
Soit $v = u - \lambda \Id_V$.
On a $v^{r-1} \neq 0$ et $v^r = 0$ [$\pi_v = X^r$].
Si $\dim \Ker v = 1$ ($=\dim \Ker(u - \lambda \Id)$).

\begin{lemma*}
$\forall k \in \intervalint{1}{r-1}$, $\dim \Ker v^{k+1} = \dim \Ker v^k + 1$
\end{lemma*}
Dém : $\forall k \in \N$, $\Ker v^k \subset \Ker v^{k+1}$.
[Car $\forall x \in V$, si $v^k(x) = 0$, $v^{k+1}(x) = v(v^k(x)) = 0$]
Si de plus $\Ker v^k = \Ker v^{k+1}$, on a toujours $\Ker v^{k+1} \subset \Ker v^{k+2}$ :
Soit $x \in \Ker(v^{k+2})$.
$v(x) \in \Ker v^{k+1} = \Ker v^k$. Donc $v^{k+1}(x) = 0$
et $\Ker v^{k+1} = \Ker v^{k+2}$.
Par récurrence immédiate :
$\forall i \geq k, \dim \Ker v^i = \dim \Ker v^k$.
De plus, soit $k \in \intervalint{1}{r-1}$ et $f_k : \Ker v^{k+1} \to \Ker v^k$, $x \mapsto v(x) \in \Lin(\Ker v^{k+1}, \Ker v^k)$.
De plus, $\Ker f_k = \Ker v$. Donc $\dim \Ker f_k = 1$.
D'après le théorème du rang, $\dim \Ker v^{k+1} = \dim \Ker f_k + \dim \Ima f_k \leq \dim \Ker v^k + 1$.
Si $\exists k \in \intervalint{0}{r-1}$, on a $\dim \Ker v^k = \dim \Ker v^{k+1}$,
Soit $p = \min \{k \in \intervalint{0}{r-1} / \Ker v^k = \Ker v^{k+1}\}$.
D'après ce qui précède, $\forall k \leq p-1$, $\Ker v^k \subsetneq \Ker v^{k+1}$.
Donc $\dim \Ker v^{k+1} = \dim \Ker v^k + 1$.
Comme $\dim \Ker v^0 = 0$, par récurrence, $\forall k \in \intervalint{0}{p}$, $\dim \Ker v^k = k$.
De plus, $\forall i \geq p$, $\Ker v^i = \Ker v^p = V$ (car $v^r = 0$, $r \geq p$, et $v^r(V) = 0 \implies \Ker v^r = V$).
Nécessairement : $p = r = n$.
Ainsi : $v^{n-1} \neq 0$ et $v^n = 0$.
Soit $x \notin \Ker v^{n-1}$.
Alors, soit $(\alpha_0, \dots, \alpha_{n-1}) \in \C^n$, si $0 = \alpha_0 x + \dots + \alpha_{n-1} v^{n-1}(x)$.
Alors soit $k = \min \{i \in \intervalint{0}{n-1} / \alpha_i \neq 0\}$ (si cet ensemble est $\neq \emptyset$).
Ainsi $\alpha_k v^k(x) + \dots + \alpha_{n-1} v^{n-1}(x) = 0$.
On applique $v^{n-1-k}$ :
$\alpha_k v^{n-1}(x) = 0$, donc $\alpha_k = 0$ : absurde.
Donc $(x, v(x), \dots, v^{n-1}(x))$ est libre, c'est donc une base de $V$.
Comme $\forall k \in \intervalint{0}{n-1}$, $u^k(x) = (v + \lambda \Id_V)^k(x) = \sum_{i=0}^k \binom{k}{i} \lambda^{k-i} v^i(x) \in \Vect(x, v(x), \dots, v^k(x))$.
La matrice de $(x, u(x), \dots, u^{n-1}(x))$ dans $(x, v(x), \dots, v^{n-1}(x))$ est de la forme
$\begin{pmatrix} 1 & & * \\ 0 & \ddots & \\ 0 & & 1 \end{pmatrix} \in \GL_n(\C)$, donc $(x, u(x), \dots, u^{n-1}(x))$ est une base de $V$.

Généralement : Si $\pi_u = (X-\lambda_1)^{r_1} \dots (X-\lambda_p)^{r_p}$
avec $(\lambda_1, \dots, \lambda_p)$ 2 à 2 distincts $\in \C^p$, $r_i \in \N^* \; \forall i \in \intervalint{1}{p}$.
$\pi_u(u) = 0$, d'après le lemme des noyaux :
$V = \bigoplus_{i=1}^p \underbrace{\Ker(u - \lambda_i \Id_V)^{r_i}}_{F_i \text{ sev stables}}$.
Soit $i \in \intervalint{1}{p}$ : $(X-\lambda_i)^{r_i}$ annule $u_i = u_{|F_i} \in \Lin(F_i)$.
Donc $\exists m \leq r_i / \pi_{u_i} = (X-\lambda_i)^m$.
Comme $\pi_u = \prod_{i=1}^p \pi_{u_i}$, nécessairement
$\forall i \in \intervalint{1}{p}$, $\pi_{u_i} = (X-\lambda_i)^{r_i}$.
Si $\forall i \in \intervalint{1}{p}$, $\dim \Ker(u - \lambda_i \Id_V) = 1$, d'après ce qui précède :
$\forall i \in \intervalint{1}{p}, \exists x_i \in F_i / (x_i, u(x_i), \dots, u^{r_i-1}(x_i))$ est une base de $F_i$.
Soit $x = \sum_{i=1}^p x_i$, soit $(\alpha_0, \dots, \alpha_{n-1}) \in \C^n$ ($\sum_{j=0}^{n-1} \alpha_j u^j(x) = 0$)
et $P = \sum_{j=0}^{n-1} \alpha_j X^j$, on a $\underbrace{P(u)(x)}_{\in V} = 0$.
$= \sum_{i=1}^p \underbrace{P(u)(x_i)}_{\in F_i \text{ par stabilité}}$.
Comme $F_1, \dots, F_p$ sont en somme directe :
$\forall i \in \intervalint{1}{p} : P(u)(x_i) = 0$.
On vérifie que $\{P \in \C[X] / P(u)(x_i) = 0\} = (X-\lambda_i)^{r_i}\C[X]$ (div euclidienne par $(X-\lambda_i)^{r_i}$).
Il s'ensuit que $\forall i \in \intervalint{1}{p}, (X-\lambda_i)^{r_i} | P$.
Donc $\underbrace{\prod_{i=1}^p (X-\lambda_i)^{r_i}}_{\text{de } d^\circ \sum_{i=1}^p r_i = n} | P$.
Ainsi $P=0$ et $(x, \dots, u^{n-1}(x))$ est libre $\to$ c'est une base de $V$.
\end{solution}

\exercise{25}{}

\begin{statementbox}
Soit $S = \R^{\N^*} = \{ (s_n)_{n \geq 1} / \forall n \in \N^*, s_n \in \R \}$
\[
\begin{aligned}
f : S &\longrightarrow S \\
(s_n)_{n \geq 1} &\longmapsto \left(\frac{s_1 + \dots + s_n}{n}\right)_{n \geq 1} \in \Lin(S).
\end{aligned}
\]
$f((s_n)_{n \geq 1}) = \left(s_1, \frac{s_1+s_2}{2}, \frac{s_1+s_2+s_3}{3}, \dots\right)$
\begin{enumerate}[label=\alph*)]
  \item Mq $f$ est bijective et expliciter, pour $u = (u_n)_{n \geq 1} \in S$, $f^{-1}(u)$.
  \item Déterminer les éléments propres de $f^{-1}$, en déduire ceux de $f$.
\end{enumerate}
\end{statementbox}

\begin{solution}
\partanswer{Q1)}
Soit $s = (s_n)_{n \geq 1} \in S$, $u = (u_n)_{n \geq 1} \in S$,
Si $f(s) = u$, alors $s_1 = u_1$, $\frac{s_1+s_2}{2} = u_2$, $\frac{s_1+\dots+s_n}{n} = u_n$.
Donc, $s_1 = u_1$, $s_2 = 2u_2 - u_1$ et $\forall n \geq 2$, $s_n = n u_n - (n-1) u_{n-1}$.
Ainsi, $f$ est injective.

Réciproquement, en posant $\forall n \geq 2$, $s_n = n u_n - (n-1) u_{n-1}$ et $s_1 = u_1$.
On a $\forall k \geq 2$, $\frac{s_1 + \dots + s_k}{k} = \frac{k u_k}{k} = u_k$.
donc $f(s) = u$.

Donc, $f$ est bijective et $f^{-1} : S \to S$,
$(u_n)_{n \geq 1} \mapsto (u_1, 2u_2 - u_1, \dots, n u_n - (n-1) u_{n-1})$.

\partanswer{Q2)}
$0 \notin \Sp(f^{-1})$ car $f^{-1}$ est bijective.
Soit $\lambda \in \R^*$, et $u \in S \setminus \{0\}$,
$f^{-1}(u) = \lambda u$ ssi $\begin{cases} u_1 = \lambda u_1 \\ \forall n \geq 2, n u_n - (n-1) u_{n-1} = \lambda u_n \end{cases}$
Définissons $k = \min \{n \geq 1 / u_n \neq 0\}$.
$k u_k = \lambda u_k$ et $\forall n \geq k+1$, $n u_n - (n-1) u_{n-1} = \lambda u_n$.
Cela implique donc $\lambda = k$ et $\forall n \geq k+1$, $u_n = \frac{n-1}{n-k} u_{n-1}$.

D'où $u_{k+1} = k u_k$.
$u_{k+2} = \frac{k+1}{2} u_{k+1} = \frac{k(k+1)}{2} u_k$.
$\forall n \in \N$, par récurrence, $u_{k+n} = \binom{k+n}{n} u_k$.

Réciproquement, en définissant pour $k \in \N^*$,
$U^{(k)}$ par : $\forall i \in \intervalint{1}{k-1}, u_i = 0$.
$u_k = 1$, $\forall n \in \N, u_{k+n} = \binom{k+n}{n}$.

Il vient $\forall n \in \N, u_{k+n+1} = \binom{k+n+1}{n+1} = \frac{k+n+1}{n+1} \binom{k+n}{n}$.
Donc, $f^{-1}(U^{(k)}) = k U^{(k)}$ en remontant les calculs précédents.
Ainsi, les valeurs propres de $f^{-1}$ sont les $k \in \N^*$ et $\forall k \in \N^*, \Ker(f^{-1} - k \Id) = \R U^{(k)}$.

De plus, $\forall u \in S \setminus \{0\}$, $\forall \lambda \in \R^*$ : $f^{-1}(u) = \lambda u$
si et seulement si $f(u) = \frac{1}{\lambda} u$.
Ainsi, $\Sp(f) = \{\frac{1}{k} / k \in \N^*\}$
et $\forall k \in \N^*$, $\Ker(f - \frac{1}{k} \Id) = \R U^{(k)}$.
\end{solution}

\exercise{26}{}

\begin{statementbox}
Soit $n \in \N^*$ et $u : \R_n[X] \to \R_n[X]$, $P \mapsto (X-1)P'$.
Déterminer les valeurs propres de $u$ et les dimensions des sous-espaces propres associés.
\end{statementbox}

\begin{solution}
On note $B$ la b.c de $\R_n[X]$.
\[
\Mat_B(u) =
\begin{pmatrix}
0 & -1 & & 0 \\
& 1 & \ddots & \\
& & \ddots & -n \\
0 & & & n
\end{pmatrix} \in \mathcal{T}_n^+(\R)
\]
Car $u(1) = 0$.
$\forall k \geq 1, u(X^k) = k X^k - k X^{k-1}$.
Il vient $\chi_u = \prod_{i=0}^n (X-i)$ donc $\Sp(u) = \intervalint{0}{n}$.
$\forall k \in \intervalint{0}{n}$, $\dim E_k = 1$ car $k$ est racine simple de $\chi_u$ ($\Ker(u - k \Id)$).

De plus, soit $P \in \R_n[X]$, $u(P) = kP$ ssi $\forall x \in \R \setminus \{1\}$ :
$(x-1)P'(x) = k P(x)$
$P'(x) = \frac{k}{x-1} P(x)$.
Donc $P$ est de la forme : $P(x) = C e^{k \ln|x-1|}$.
Donc $P(x) = C (x-1)^k$.
Soit $P_k = (X-1)^k$.
$u(P_k) = k P_k$.
Donc $\Ker(u - k \Id) = \R P_k$.
\end{solution}

\exercise{43}{}

\begin{statementbox}
Soit $A = \begin{pmatrix} 2 & 1 & 0 \\ -3 & -1 & -1 \\ 1 & 0 & -1 \end{pmatrix}$ et $u \in \Lin(\R^3)$ c.a à $A$.
\begin{enumerate}[label=\alph*)]
  \item Trouver tous les sev de $\R^3$ stables par $u$ [On pourra calculer $\chi_u$, et vérifier que $A^3 = 0, A^2 \neq 0$].
  \item Existe-t-il $M \in \M_3(\R) / M^2 = A$ ?
\end{enumerate}
\end{statementbox}

\begin{solution}
\begin{align*}
\chi_A = \det(X I_3 - A) &=
\begin{vmatrix} X-2 & -1 & 0 \\ 3 & X+1 & 1 \\ -1 & 0 & X+1 \end{vmatrix} \\
&= \begin{vmatrix} X & X & X \\ 3 & X+1 & 1 \\ -1 & 0 & X+1 \end{vmatrix} \quad L_1 \leftarrow L_1 + L_2 + L_3 \\
&= X \begin{vmatrix} 1 & 1 & 1 \\ 3 & X+1 & 1 \\ -1 & 0 & X+1 \end{vmatrix}
\end{align*}
Ainsi :
\begin{align*}
\chi_A &= X \begin{vmatrix} 1 & 0 & 0 \\ 3 & X-2 & -2 \\ -1 & 1 & X+2 \end{vmatrix} \quad \begin{matrix} C_2 \leftarrow C_2 - C_1 \\ C_3 \leftarrow C_3 - C_1 \end{matrix} \\
&= X (X^2 - 4 + 4) \\
&= X^3
\end{align*}
Il s'ensuit :
$\Sp(u) = \{0\}$.
D'après T.C.H (cf cours) $A^3 = 0$.

De plus, soit $X = \begin{pmatrix} x \\ y \\ z \end{pmatrix} \in \M_{3,1}(\R)$.
$AX = \begin{pmatrix} 0 \\ 0 \\ 0 \end{pmatrix}$ ssi $\begin{cases} 2x + y = 0 \\ -3x - y - z = 0 \\ x - z = 0 \end{cases}$
ssi $\begin{cases} x = z \\ y = -2x \end{cases}$
ssi $X \in \Vect \begin{pmatrix} 1 \\ -2 \\ 1 \end{pmatrix}$, $\dim \Ker u = 1$.

$A^2 = \begin{pmatrix} 2 & 1 & 0 \\ -3 & -1 & -1 \\ 1 & 0 & -1 \end{pmatrix} \begin{pmatrix} 2 & 1 & 0 \\ -3 & -1 & -1 \\ 1 & 0 & -1 \end{pmatrix} = \begin{pmatrix} 1 & 1 & -1 \\ -4 & -2 & 2 \\ 1 & 1 & -1 \end{pmatrix} \neq 0$ de rang $1$.
Ainsi $\dim \Ker u = 1$, $\dim \Ker u^2 = 2$, $\dim \Ker u^3 = 3$.

* Soit $F$ un sev stable par $u$.
\begin{itemize}
  \item Si $\dim F = 0$, $F = \{0\}$.
  \item Si $\dim F = 1$, $F$ est engendré par un vecteur propre. $F = \Ker u$.
  \item Si $\dim F = 2$. Comme $u$ est nilpotent, $u_{|F}$ l'est aussi, et $\dim \Ker u_{|F} \geq 1$.
  On a vu que : $\dim \Ker u_{|F}^2 = 2 = \dim F$.
  Donc $u_{|F}^2 = 0$.
  $F \subset \Ker u^2$. Donc $F = \Ker u^2$.
  \item Si $\dim F = 3$, $F = \R^3$.
\end{itemize}

\partanswer{b)}
S'il existe $M \in \M_3(\R) / A = M^2$.
Comme $A^3 = 0$, $M^6 = 0$.
$M$ est nilpotente, on verra (cf cours) que $M^3 = 0$.
Alors $M^4 = A^2 = 0 \to$ impossible.
Ainsi il n'existe pas de racine carrée ($A = M^2$).
\end{solution}

\exercise{44}{}

\begin{statementbox}
Soit $M = \begin{pmatrix} 1 & -1 & 0 \\ -1 & 2 & 1 \\ 1 & 0 & 1 \end{pmatrix}$ et $u \in \Lin(\R^3)$ c.a. à $M$.
a) Trouver tous les sev de $\R^3$ stables par $u$ [On pourra évaluer $\chi_M$]
\end{statementbox}

\begin{solution}
\begin{align*}
\det(X I_3 - M) &= \begin{vmatrix} X-1 & 1 & 0 \\ 1 & X-2 & -1 \\ -1 & 0 & X-1 \end{vmatrix} \\
&= \begin{vmatrix} X-1 & 1 & 0 \\ 1 & X-2 & -1 \\ 0 & X-2 & X-2 \end{vmatrix} \quad L_3 \leftarrow L_3 + L_2 \\
&= (X-2) \begin{vmatrix} X-1 & 1 & 0 \\ 1 & X-2 & -1 \\ 0 & 1 & 1 \end{vmatrix} \\
&= (X-2) \begin{vmatrix} X-1 & 1 & -1 \\ 1 & X-2 & X-1 \\ 0 & 1 & 0 \end{vmatrix} \quad C_3 \leftarrow C_3 - C_2 \\
&= -(X-2) \begin{vmatrix} X-1 & -1 \\ 1 & X+1 \end{vmatrix} \\
&= (X-2) (X^2 - 2X) \\
&= X (X-2)^2
\end{align*}
D'après le théorème de Cayley-Hamilton et le lemme des noyaux,
\[
\R^3 = \underbrace{\Ker(u)}_{F_0} \oplus \underbrace{\Ker((u-2\Id)^2)}_{F_2}
\]
deux sev stables par $u$.
\end{solution}

\exercise{56}{}

\begin{statementbox}
Soit $n \in \N^*$, $(x_1, \dots, x_n) \in \C^n$, $A = \begin{pmatrix} 0 & & & x_n \\ & & \iddots & \\ & \iddots & & \\ x_1 & & & 0 \end{pmatrix}$
% À vérifier sur la photo : la donnée (x_1, \dots, x_{2n}) \in \C^n a été corrigée en x_n pour avoir la bonne dimension, comme utilisé dans la suite de la correction (A_i de taille 2 et x_i x_{2p+1-i}).

Déterminer $\Sp_\C A$.
\end{statementbox}

\begin{solution}
Si $n$ est pair, $n = 2p$.
Soit $u \in \Lin(\C^n)$ c.a. à $A$ et $(e_1, \dots, e_n)$ la b.c de $\C^n$.
$u(e_1) = x_1 e_n, \dots, u(e_n) = x_n e_1$.
\begin{align*}
\text{Soit } & F_1 = \Vect(e_1, e_n) \\
& F_2 = \Vect(e_2, e_{n-1}) \\
& \vdots \\
& F_p = \Vect(e_p, e_{p+1})
\end{align*}
Stables par $u$.
On a : $\C^n = \bigoplus_{i=1}^p F_i$.
et $\Mat_{(e_i, e_{2p+1-i})}(u_{|F_i}) = \begin{pmatrix} 0 & x_{2p+1-i} \\ x_i & 0 \end{pmatrix} = A_i$.
\[
\chi_{A_i} = \begin{vmatrix} X & -x_{2p+1-i} \\ -x_i & X \end{vmatrix} = X^2 - x_i x_{2p+1-i}
\]
Si $x_i$ ou $x_{2p+1-i} = 0$, $\Sp(u_{|F_i}) = \{0\}$.
Si $x_i \neq 0$ et $x_{2p+1-i} \neq 0$ :
\[
\Sp(u_{|F_i}) = \{\lambda_i, -\lambda_i\} \text{ avec } \lambda_i \text{ une racine carrée sur } \C \text{ de } x_i x_{2p+1-i}.
\]

Si $n$ est impair, $n = 2p + 1$.
De même, $F_1 = \Vect(e_1, e_{2p+1})$
$F_2 = \Vect(e_2, e_{2p})$
$\vdots$
$F_p = \Vect(e_p, e_{p+2})$
$F_{p+1} = \Vect(e_{p+1})$
Les valeurs propres sont les mêmes que dans le cas pair, en ajoutant $x_{p+1}$.
Dans tous les cas $\chi_u = \prod_{i=1}^p \chi_{u_{|F_i}}$.
\end{solution}
